\subsection{命题逻辑的推理理论}

\subsubsection{定义}

设 $A$ 和 $B$ 是两个命题公式，当且仅当 $A \to B$ 是重言式时，称从 $A$ 可推出 $B$，或 $B$ 是前提 $A$ 的结论，记为 $A \Rightarrow B$。

\textbf{证由前提 $A$ 推出结论 $B$ 的方法:}

\begin{enumerate}
	\item 真值表法；
	\item 等值演算法；
	\item 在自然推理系统 $P$ 中用推理规则证明。
\end{enumerate}

\subsubsection{推理规则}

\begin{enumerate}
	\item $p \land q \Rightarrow p,\quad p \land q \Rightarrow q$ \hfill （化简规则）
	\item $p \Rightarrow p \lor q,\quad q \Rightarrow p \lor q$ \hfill （附加规则）
	\item $p,\ p \to q \Rightarrow q$ \hfill （假言推理）
	\item $p \to q,\ \neg q \Rightarrow \neg p$ \hfill （拒取式）
	\item $p \lor q,\ \neg p \Rightarrow q$ \hfill （析取三段论）
	\item $p,\ q \Rightarrow p \land q$ \hfill （合取式）
	\item $p \to q,\ q \to r \Rightarrow p \to r$ \hfill （假言三段论）
	\item $p \leftrightarrow q,\ q \leftrightarrow r \Rightarrow p \leftrightarrow r$ \hfill （等价三段论）
	\item $p \to q,\ r \to s,\ p \lor r \Rightarrow q \lor s$ \hfill （构造性二难）
	\item $p \lor q,\ \neg p \lor s \Rightarrow q \lor s$ \hfill （归结式）
\end{enumerate}

\begin{question}{例题1}
证明：\((p \vee q) \land (p \leftrightarrow r) \land (q \rightarrow s) \Rightarrow s \vee r\)

\textbf{解答：}

令
\[
A=(p \vee q) \land (p \leftrightarrow r) \land (q \rightarrow s),
\qquad
B=s\vee r
\]
要证明
\[
A\Rightarrow B
\]
只需证明
\[
A\land \neg B
\]
为矛盾式。

即证明：
\[
(p \vee q) \land (p \leftrightarrow r) \land (q \rightarrow s)\land \neg(s\vee r)
\]
为矛盾式。

下面进行等值演算：
\begin{align*}
	& (p \vee q) \land (p \leftrightarrow r) \land (q \rightarrow s)\land \neg(s\vee r) \\
	\Leftrightarrow{}&
	(p \vee q) \land \bigl((p\rightarrow r)\land(r\rightarrow p)\bigr)
	\land (\neg q\vee s)\land(\neg s\land \neg r) \\
	\Leftrightarrow{}&
	(p \vee q) \land (\neg p\vee r)\land(\neg r\vee p)
	\land(\neg q\vee s)\land\neg s\land \neg r \\
	\Leftrightarrow{}&
	(p \vee q) \land \bigl((\neg p\vee r)\land \neg r\bigr)
	\land(\neg r\vee p)
	\land\bigl((\neg q\vee s)\land\neg s\bigr) \\
	\Leftrightarrow{}&
	(p \vee q) \land (\neg p\land \neg r)
	\land(\neg r\vee p)
	\land(\neg q\land \neg s) \\
	\Leftrightarrow{}&
	(p \vee q)\land \neg p\land \neg q\land \neg r\land \neg s \\
	\Leftrightarrow{}&
	\bigl((p\vee q)\land \neg p\land \neg q\bigr)
	\land \neg r\land \neg s \\
	\Leftrightarrow{}&
	\bigl((p\land \neg p\land \neg q)\vee(q\land \neg p\land \neg q)\bigr)
	\land \neg r\land \neg s \\
	\Leftrightarrow{}&
	0
\end{align*}
因此，
\[
(p \vee q) \land (p \leftrightarrow r) \land (q \rightarrow s)\land \neg(s\vee r)
\]
为矛盾式。
所以，
\[
(p \vee q) \land (p \leftrightarrow r) \land (q \rightarrow s) \Rightarrow s \vee r
\]
解释：
如果结论 \(s\vee r\) 为假，则有
\[
\neg(s\vee r)\Leftrightarrow \neg s\land \neg r
\]
由
\[
p\leftrightarrow r
\]
可得
\[
p\rightarrow r
\]
即
\[
\neg p\vee r
\]
再结合 \(\neg r\)，可推出
\[
\neg p
\]
由
\[
q\rightarrow s
\]
即
\[
\neg q\vee s
\]
再结合 \(\neg s\)，可推出
\[
\neg q
\]
于是得到
\[
\neg p\land \neg q
\]
但前提中又有
\[
p\vee q
\]
这与
\[
\neg p\land \neg q
\]
矛盾。

因此当前提
\[
(p \vee q) \land (p \leftrightarrow r) \land (q \rightarrow s)
\]
成立时，结论
\[
s\vee r
\]
必成立。


\end{question}

\begin{question}{例题2}
给出下面推论的形式证明：

若马会飞或羊吃草，则母鸡就会变飞鸟；

如果母鸡变了飞鸟，那么烤熟的鸭子还会跑；

烤熟的鸭子不会跑，所以羊儿不吃草。

\textbf{解答：}

设命题变元如下：
\[
p:\text{马会飞}
\]
\[
q:\text{羊吃草}
\]
\[
r:\text{母鸡变飞鸟}
\]
\[
s:\text{烤熟的鸭子会跑}
\]

则题中的推论可符号化为：
\[
(p\vee q)\to r,\quad r\to s,\quad \neg s \vdash \neg q
\]

即需要证明：
\[
\bigl((p\vee q)\to r\bigr)\land(r\to s)\land \neg s \Rightarrow \neg q
\]

形式证明如下：
\[
\begin{array}{c|l|l}
	序号 & 公式 & 理由 \\
	\hline
	1 & (p\vee q)\to r & 前提 \\
	2 & r\to s & 前提 \\
	3 & \neg s & 前提 \\
	4 & \neg r & 由 2,3 用拒取式 \\
	5 & \neg(p\vee q) & 由 1,4 用拒取式 \\
	6 & \neg p\land \neg q & 由 5 用德摩根律 \\
	7 & \neg q & 由 6 用化简律
\end{array}
\]

所以可得：
\[
\neg q
\]

因此原推论成立，即：
\[
(p\vee q)\to r,\quad r\to s,\quad \neg s \vdash \neg q
\]

用文字解释就是：

由
\[
r\to s
\]
和
\[
\neg s
\]
可推出
\[
\neg r
\]
即母鸡没有变飞鸟。

又由
\[
(p\vee q)\to r
\]
和
\[
\neg r
\]
可推出
\[
\neg(p\vee q)
\]

根据德摩根律：
\[
\neg(p\vee q)\Leftrightarrow \neg p\land \neg q
\]

所以得到：
\[
\neg q
\]

即羊儿不吃草。


\end{question}

\begin{question}{例题3}

用附加前提证明法证明：

如果小张和小王去看电影，则小李也去看电影；

小赵不去看电影或小张去看电影，小王去看电影；

所以当小赵去看电影时，小李也去。

\textbf{解答：}

设命题变元如下：
\[
p:\text{小张去看电影}
\]
\[
q:\text{小王去看电影}
\]
\[
r:\text{小李去看电影}
\]
\[
s:\text{小赵去看电影}
\]
则题中前提和结论可符号化为：
\[
(p\land q)\to r
\]
\[
\neg s\vee p
\]
\[
q
\]
结论为：
\[
s\to r
\]
即需要证明：
\[
(p\land q)\to r,\quad \neg s\vee p,\quad q \vdash s\to r
\]
使用附加前提证明法。
为了证明
\[
s\to r
\]
只需附加前提
\[
s
\]
然后推出
\[
r
\]
形式证明如下：
\[
\begin{array}{c|l|l}
	序号 & 公式 & 理由 \\
	\hline
	1 & (p\land q)\to r & 前提 \\
	2 & \neg s\vee p & 前提 \\
	3 & q & 前提 \\
	4 & s & 附加前提 \\
	5 & p & 由 2,4 用析取三段论 \\
	6 & p\land q & 由 5,3 用合取引入律 \\
	7 & r & 由 1,6 用假言推理 \\
	8 & s\to r & 由 4--7 用附加前提证明法
\end{array}
\]
因此，
\[
(p\land q)\to r,\quad \neg s\vee p,\quad q \vdash s\to r
\]
所以原推论成立。
文字解释：
为了证明“当小赵去看电影时，小李也去”，即证明
\[
s\to r
\]
我们先假设小赵去看电影，即
\[
s
\]
又因为前提中有
\[
\neg s\vee p
\]
即“小赵不去看电影或小张去看电影”。
在已经假设
\[
s
\]
成立的情况下，\(\neg s\) 为假，因此只能推出
\[
p
\]
即小张去看电影。
又已知
\[
q
\]
即小王去看电影，所以有
\[
p\land q
\]
由前提
\[
(p\land q)\to r
\]
可推出
\[
r
\]
即小李去看电影。
所以由附加前提证明法可得：
\[
s\to r
\]
即当小赵去看电影时，小李也去看电影。


\end{question}
